\documentclass[a4paper,12pt]{article}
\usepackage{docsTemplate}
\usepackage{longtable}

\setTitle{Glossario}
\setAuthors{Filippo Compagno \\ & Mattia Oliva Medin \\ & Sara Gioia Fichera \\ Francesco Balestro}
\setVerificators{Andrea Masiero \\ & Nenad Radulovic \\ & Bianca Zaghetto}
\setApprovation{}
\setVersion{0.5}
\setType{Interno}
\setDestination{Team Three Technologies \\ & Tullio Vardanega \\ & Riccardo Cardin}

\begin{document}
    \include{docTitlepage}

        \newpage
        \section*{Registro delle modifiche} {
        \begin{table}[h!]
            \rowcolors{2}{lightgray}{white}
            \begin{tabularx}{\textwidth}{|l|l|l|l|L|L|}
                \hline
                \textbf{Vers.} & \textbf{Data} & \textbf{Autore} & \textbf{Ruolo} & \textbf{Verifica} & \textbf{Oggetto} \\
                \hline
                \textbf{0.5} & 17/03/2026 & Francesco Balestro & Amministratore & Andrea Masiero & Aggiunta di altri termini \\
                \hline
                \textbf{0.4} & 21/01/2026 & Sara Gioia Fichera & Amministratore & Bianca Zaghetto & Aggiunta termini relativi alle tecnologie \\
                \hline
                \textbf{0.3} & 17/12/2025 & Mattia Oliva Medin & Amministratore & Nenad Radulovic & Aggiunta termini relativi al prodotto \\
                \hline
                \textbf{0.2} & 08/11/2025 & Filippo Compagno & Amministratore & Andrea Masiero & Aggiunta di altri termini \\
                \hline
                \textbf{0.1} & 04/11/2025 & Filippo Compagno & Amministratore & Andrea Masiero & Creazione della struttura e prima stesura \\
                \hline
            \end{tabularx}
            \end{table}
        }
        \newpage

        \tableofcontents
        \newpage

        \section*{Introduzione} {
            Questo documento ha lo scopo di chiarire il significato di ciascun termine nel dominio del progetto che può risultare ambiguo o di difficile comprensione.\\
            All'interno degli altri documenti prodotti la notazione
            \begin{center}
                parola\textsubscript{G}
            \end{center}
            indica la presenza della parola nel Glossario.\\
            Il documento è organizzato per sezioni, una per ogni lettera dell'alfabeto, ciascuna delle quali riporta una lista dei termini ordinati in ordine alfabetico.
        }
        \addcontentsline{toc}{section}{\numberline{}Introduzione}
        \newpage

        \section*{A} {
            \begin{itemize}
                \item \textbf{AdR}\\
                    Acronimo di Analisi dei Requisiti, per la sua definizione vedasi "Analisi dei Requisiti".
                \item \textbf{Agile}\\
                    Insieme di metodi di sviluppo software in contrapposizione ai modelli di sviluppo tradizionali, basati su pianificazione adattiva, coinvolgimento continuo degli stakeholder e sviluppo iterativo e incrementale.
                \item \textbf{Amministratore}\\
                    Figura professionale che si occupa di mantenere procedure, strumenti e tecnologie a supporto del lavoro del team durante tutto lo svolgimento del progetto.
                \item \textbf{Angular}\\
                    Framework open-source sviluppato da Google in inguaggio TypeScript, utilizzato per lo sviluppo di applicazioni web.
                \item \textbf{Analisi dei Requisiti}\\
                    Documento redatto dal team di Analisti che contiene tutti i casi d'uso relativi al capitolato e individuati dagli analisti.
                \item \textbf{Analista}\\
                    Figura professionale che si occupa della raccolta, analisi e definizione dei requisiti del progetto.
                \item \textbf{Archivio}\\
                    Complesso dei documenti prodotti o acquisiti da un soggetto pubblico o privato durante lo svolgimento della propria attività.
                \item \textbf{Attore}\\
                    Entità che interagisce con il sistema in modo attivo, per soddisfare un suo bisogno (attore principale), o in modo passivo, invocato dal sistema per fornire supporto (attore secondario).
            \end{itemize}
        }
        \addcontentsline{toc}{section}{\numberline{}A}
        \newpage

        \section*{B} {
            \begin{itemize}
                \item \textbf{Baseline}\\
                    Versione approvata di un prodotto di progetto che può essere modificata solo tramite metodi stabiliti dal sistema di controllo delle modifiche.
            \end{itemize}
        }
        \addcontentsline{toc}{section}{\numberline{}B}
        \newpage

        \section*{C} {
            \begin{itemize}
                \item \textbf{Capitolato}\\
                    Documento redatto dal committente di un appalto o di un progetto in cui vengono fissate le aspettative, le specifiche, i requisiti e i vincoli sul prodotto da realizzare.
                \item \textbf{Caso d'uso}\\
                    Insieme di scenari, cioè una sequenza di azioni, che hanno in comune uno scopo finale (obiettivo) per un utente (attore). In altre parole è una situazione nella quale il sistema viene utilizzato per soddisfare i bisogni dell'utente.
                \item \textbf{Committente}\\
                    Soggetto che richiede lo sviluppo di un software e ne autorizza (e finanzia) la realizzazione.
                \item \textbf{Conservazione digitale}\\
                    Insieme delle attività finalizzate a definire ed attuare le politiche complessive del sistema di conservazione digitale e a governarne la gestione in relazione al modello organizzativo adottato, garantendo nel tempo le caratteristiche di autenticità, integrità, leggibilità, reperibilità dei documenti.
                \item \textbf{Consuntivo}\\
                    Resoconto di fine sprint che riassume risultati, attività svolte, tempi, costi e scostamenti rispetto a quanto pianificato, utilizzato per valutare l'andamento effettivo del progetto e del periodo di lavoro.
                \item \textbf{Cruscotto}\\
                    Interfaccia grafica che mostra l'andamento delle metriche durante il corso del progetto tramite strumenti come tabelle e grafici.
            \end{itemize}
        }
        \addcontentsline{toc}{section}{\numberline{}C}
        \newpage

        \section*{D} {
            \begin{itemize}
                \item \textbf{Dashboard}\\
                    Inglese di Cruscotto, per la sua definizione vedasi "Cruscotto".
                \item \textbf{Digest}\\
                    Stringa di lunghezza fissa prodotta da una funzione di hash.
                \item \textbf{DIP}\\
                    Acronimo di Distribution Information Package, per la sua definizione vedasi "Distribution Information Package".
                \item \textbf{Discord}\\
                    Piattaforma di comunicazione online che offre funzionalità di chat testuale, chat vocale e video chiamate. La comunicazione può avvenire in forma privata o in comunità virtuali chiamate "server", accessibili tramite link di invito.
                \item \textbf{Distribution Information Package}\\
                    Pacchetto di distribuzione prodotto dal sistema di conservazione in formato di archivio \verb|.zip|. Contiene tipicamente documenti, metadati, report e altre informazioni accessorie.
            \end{itemize}
        }
        \addcontentsline{toc}{section}{\numberline{}D}
        \newpage

        \section*{E} {
            \begin{itemize}
                \item \textbf{Efficacia}\\
                    Grado con cui un'azione, un processo o un sistema raggiunge gli obiettivi prefissati o produce i risultati attesi, indipendentemente dalle risorse impiegate.
                \item \textbf{Efficienza}\\
                    Capacità di ottenere un risultato impiegando il minimo di risorse possibili, ottimizzando il rapporto tra input utilizzati e output prodotti.
                \item \textbf{Electron}\\
                    Framework open-source gestito e ospitato da GitHub, consente lo sviluppo della GUI di applicazioni desktop utilizzando tecnologie web tra cui motore di rendering Chromium e il runtime Node.js.
            \end{itemize}
        }
        \addcontentsline{toc}{section}{\numberline{}E}
        \newpage

        \section*{F} {
            \begin{itemize}
                \item \textbf{File}\\
                    Insieme di informazioni, dati o comandi logicamente correlati, raccolti sotto un unico nome e registrati, per mezzo di un programma di elaborazione o discrittura, nella memoria di un computer.
                \item \textbf{Fornitore}\\
                    Soggetto che progetta, sviluppa e consegna il software richiesto.
            \end{itemize}
        }
        \addcontentsline{toc}{section}{\numberline{}F}
        \newpage

        \section*{G} {
            \begin{itemize}
                \item \textbf{Gestione Documentale}\\
                    Processo finalizzato al controllo efficiente e sistematico della produzione, ricezione, tenuta, uso, selezione e conservazione dei documenti.
                \item \textbf{GitHub}\\
                    Piattaforma di hosting per repository Git che rende semplice la collaborazione tra sviluppatori e la gestione di progetti in maniera organizzata.
            \end{itemize}
        }
        \addcontentsline{toc}{section}{\numberline{}G}
        \newpage

        \section*{H} {
            \begin{itemize}
                \item \textbf{Hash}\\
                    Funzione non invertibile che mappa una stringa di lunghezza arbitraria in una stringa di lunghezza predefinita.
                \item \textbf{Hashing}\\
                    Processo in cui si applica una funzione di hash.
            \end{itemize}
        }
        \addcontentsline{toc}{section}{\numberline{}H}
        \newpage

        \section*{I} {
            \begin{itemize}
                \item \textbf{Identificativo univoco}\\
                    Sequenza di numeri o caratteri alfanumerici associata in modo univoco e persistente ad un'entità all'interno di uno specifico ambito di applicazione. 
            \end{itemize}            
        }
        \addcontentsline{toc}{section}{\numberline{}I}
        \newpage

        \section*{J} {
            
        }
        \addcontentsline{toc}{section}{\numberline{}J}
        \newpage

        \section*{K} {
            
        }
        \addcontentsline{toc}{section}{\numberline{}K}
        \newpage

        \section*{L} {
            \begin{itemize}
                \item \textbf{LaTeX}\\
                    Linguaggio di markup per la preparazione di testi, basato sul programma di composizione tipografica TeX.
            \end{itemize}
        }
        \addcontentsline{toc}{section}{\numberline{}L}
        \newpage

        \section*{M} {
            \begin{itemize}
                \item \textbf{Metadati}\\
                    Dati associati a un documento informatico, a un fascicolo informatico o a un'aggregazione documentale per identificarli, descrivendone il contesto, il contenuto e la struttura - così da permetterne la gestione del tempo - in conformità a quanto definito nella norma ISO 15489-1:2016 e più nello specifico dalla norma ISO 23081-1:2017.
                \item \textbf{Milestone}\\
                    Data di calendario che fissa un punto di avanzamento atteso nel tempo di progetto.
                \item \textbf{Minimum Viable Product}\\
                    Versione di un prodotto software con caratteristiche appena sufficienti per essere utilizzabile dai primi clienti, i quali possono fornire feedback per lo sviluppo futuro.
                \item \textbf{MVP}\\
                    Acronimo di Minimum Viable Product, per la sua definizione vedasi "Minimum Viable Product".
            \end{itemize}
        }
        \addcontentsline{toc}{section}{\numberline{}M}
        \newpage

        \section*{N} {
            
        }
        \addcontentsline{toc}{section}{\numberline{}N}
        \newpage

        \section*{O} {
            
        }
        \addcontentsline{toc}{section}{\numberline{}O}
        \newpage

        \section*{P} {
            \begin{itemize}
                \item \textbf{Pacchetto di distribuzione}\\
                    Pacchetto informativo inviato dal sistema di conservazione all’utente in risposta ad una sua richiesta di accesso a oggetti di conservazione.
                \item \textbf{PoC}\\
                    Acronimo di Proof of Concept, per la sua definizione vedasi "Proof of Concept".
                \item \textbf{Progettista}\\
                    Figura professionale che si occupa di progettare l'architettura del sistema a partire dall'analisi effettuata.
                \item \textbf{Programmatore}\\
                    Figura professionale che si occupa della codifica del prodotto software obiettivo del progetto, realizzando l'architettura progettata.
                \item \textbf{Proponente}\\
                    Soggetto che propone la realizzazione di un progetto, tramite capitolato, di suo interesse.
                \item \textbf{Proof of Concept}\\
                    Metodo per dimostrare la fattibilità o confermare la validità di alcuni principi o concetti fondamentali, si riferisce alla dimostrazioni pratica dei funzionamenti di base di un applicativo software o di un intero sistema.
            \end{itemize}
        }
        \addcontentsline{toc}{section}{\numberline{}P}
        \newpage

        \section*{Q} {
            
        }
        \addcontentsline{toc}{section}{\numberline{}Q}
        \newpage

        \section*{R} {
            \begin{itemize}
                \item \textbf{Repository}\\
                    Archivio digitale che contiene file e l'intera cronologia delle modifiche associate. Include dati e metadati necessari al sistema di controllo versione per tracciare revisioni, collaborazioni e gestione dei rami di sviluppo.
                \item \textbf{Requirements and Technology Baseline}\\
                    Prima revisione di avanzamento del progetto didattico di IS. Ad essa sono associate due baseline, una riguardante i requisiti del prodotto e l'altra le tecnologie scelte, la cui fattibilità di utilizzo è dimostrata tramite Proof of Concept.
                \item \textbf{Requisito}\\
                    Capacità, dal lato del bisogno o della soluzione, che definisce ciò che un utente deve poter ottenere e ciò che un sistema deve rispettare per soddisfare obiettivi, specifiche o vincoli formali.
                \item \textbf{Responsabile}\\
                    Figura professionale che si occupa della gestione e il coordinamento delle attività del team e delle comunicazione verso l'esterno.
                \item \textbf{RTB}\\
                    Acronimo di Requirements and Technology Baseline, per la sua definizione vedasi "Requirements and Technology Baseline".
                \item \textbf{Retrospettiva}\\
                    Momento di analisi e riflessione a posteriori su processi e modalità operative, finalizzato a individuare problemi, punti di forza e azioni di miglioramento per le attività future.
            \end{itemize}
        }
        \addcontentsline{toc}{section}{\numberline{}R}
        \newpage

        \section*{S} {
            \begin{itemize}
                \item \textbf{Sistema di conservazione}\\
                    Insieme di regole, procedure e tecnologie che assicurano la conservazione dei documenti informatici in attuazione a quanto previsto dall’art. 44, comma 1, del CAD.
                \item \textbf{Sprint}\\
                    Breve periodo di tempo, tipicamente tra una e quattro settimane, in cui il gruppo completa un insieme di attività pianificate al suo inizio.
                \item \textbf{Stakeholder}\\
                    Persona, gruppo di persone o organizzazione che ha interesse e influenza sul progetto e il prodotto.
            \end{itemize}
        }
        \addcontentsline{toc}{section}{\numberline{}S}
        \newpage

        \section*{T} {
            \begin{itemize}
                \item \textbf{Telegram}\\
                    Piattaforma di messaggistica istantanea, utilizzata dal gruppo per scambiare file in modo organizzato, tramite un gruppo con chat divise in sezioni.
                \item \textbf{Test}\\
                    Esecuzione di prova di un sistema software o di una sua parte volta a trovare difetti nel codice.
            \end{itemize}
        }
        \addcontentsline{toc}{section}{\numberline{}T}
        \newpage

        \section*{U} {
            \begin{itemize}
                \item \textbf{UML}\\
                    Acronimo di Unified Modeling Language, per la sua definizione vedasi "Unified Modeling Language".
                \item \textbf{Unified Modeling Language}\\
                    Linguaggio di modellazione e specifica, basato sul paradigma ad oggetti, che definisce diversi diagrammi tra cui quello delle classi e dei casi d'uso.

            \end{itemize}
        }
        \addcontentsline{toc}{section}{\numberline{}U}
        \newpage

        \section*{V} {
            \begin{itemize}
                \item \textbf{Validazione}\\
                    Controllo atto a garantire che quanto sviluppato soddisfi i requisiti.
                \item \textbf{Verifica}\\
                    Controllo atto a garantire che un processo o un prodotto soddisfi determinati standard.
                \item \textbf{Verificatore}\\
                    Figura professionale che si occupa di revisionare documenti e/o codice sorgente prodotto, segnalando tutti gli errori in modo che questi vengano corretti per perseguire un alto standard di qualità e il successo del progetto.
                \item \textbf{Versionamento}\\
                    Processo di gestione e tracciamento delle modifiche ai prodotti del progetto, per mantenere una cronologia dei cambiamenti e consentire il ripristino di stati precedenti.
            \end{itemize}
        }
        \addcontentsline{toc}{section}{\numberline{}V}
        \newpage

        \section*{W} {
            \begin{itemize}
                \item \textbf{Way of working}\\
                    Insieme delle metodologie di lavoro, dei processi e degli strumenti che viene adottato da un gruppo per svolgere il proprio compito.
                \item \textbf{WhatsApp}\\
                    Piattaforma di messaggistica istantanea, utilizzata dal team tramite un gruppo per messaggi veloci.
            \end{itemize}
        }
        \addcontentsline{toc}{section}{\numberline{}W}
        \newpage

        \section*{X} {
            
        }
        \addcontentsline{toc}{section}{\numberline{}X}
        \newpage

        \section*{Y} {
            
        }
        \addcontentsline{toc}{section}{\numberline{}Y}
        \newpage

        \section*{Z} {
            
        }
        \addcontentsline{toc}{section}{\numberline{}Z}

\end{document}